% ============================================================================
% Section 2: Hardware SPI Configuration (Task 2)
% Handout s10.2 requires ONLY: the Task 2 configuration table, the SPI clock
% calculation, the relevant alternate-function information, and the manual
% sections or datasheet tables used. README s5.1 also asks to record the
% continuity check of the EEPROM pins.
% ============================================================================
\section{Hardware SPI Configuration}
\label{sec:spi-config}

\paragraph{Pin mapping.}
The EEPROM connections given in the README~\cite{readme} were confirmed by
continuity (board unpowered, P2 jumpers removed), after anchoring on EEPROM
VCC and VSS (\cref{tab:pins}).

\begin{table}[htbp]
  \centering
  \caption{EEPROM to STM32 connections, verified by continuity.}
  \label{tab:pins}
  \begin{tabular}{@{}lllll@{}}
    \toprule
    \textbf{EEPROM pin} & \textbf{Function} & \textbf{STM32 pin} & \textbf{SPI signal} & \textbf{Continuity} \\
    \midrule
    1 & CS\#                    & PB12     & GPIO chip select & beep \\
    6 & SCK                     & PB13     & SPI2\_SCK        & beep \\
    2 & SO (data out of EEPROM) & PB14     & SPI2\_MISO       & beep \\
    5 & SI (data into EEPROM)   & PB15     & SPI2\_MOSI       & beep \\
    \midrule
    8 & VCC                     & 3V3 rail & --               & beep \\
    4 & VSS                     & GND      & --               & beep \\
    \bottomrule
  \end{tabular}
\end{table}

\begin{table}[htbp]
  \centering
  \caption{SPI configuration.}
  \label{tab:spi-config}
  \small
  \begin{tabularx}{\linewidth}{@{}l l X >{\raggedright\arraybackslash}p{3.1cm}@{}}
    \toprule
    \textbf{Item} & \textbf{Value} & \textbf{Justification} & \textbf{Source} \\
    \midrule
    SPI peripheral instance    & \SPIInstance & Only SPI peripheral on PB13--PB15
      & \cite[Tab.~15]{stm32f051ds} \\
    Peripheral clock frequency & \PCLKMHz~MHz & HSI, AHB $\div$1, APB $\div$1; SPI2 on APB (PCLK)
      & \cite[\S6.2, \S6.4.2, Tab.~1]{rm0091} \\
    PB13 alternate function    & \SPIAF & SPI2\_SCK   & \cite[Tab.~15]{stm32f051ds} \\
    PB14 alternate function    & \SPIAF & SPI2\_MISO  & \cite[Tab.~15]{stm32f051ds} \\
    PB15 alternate function    & \SPIAF & SPI2\_MOSI  & \cite[Tab.~15]{stm32f051ds} \\
    Clock divider              & $\div$\BRDivider{} (BR = \texttt{\BRBits}) & \Cref{eq:sck}
      & \cite[\S28.9.1]{rm0091} \\
    Predicted SCK frequency    & \fSCKPred~kHz & \Cref{eq:sck} & -- \\
    CPOL                       & 0 & EEPROM supports modes (0,0) and (1,1)
      & \cite[p.~5]{eeprom}; \cite[\S28.9.1]{rm0091} \\
    CPHA                       & 0 & EEPROM latches SI on the rising (first) edge
      & \cite[p.~5]{eeprom}; \cite[\S28.9.1]{rm0091} \\
    Bit order                  & MSB first & EEPROM transfers MSB first
      & \cite[Figs.~3, 5]{eeprom}; \cite[\S28.9.1]{rm0091} \\
    \bottomrule
  \end{tabularx}
\end{table}

\paragraph{SPI clock calculation.}
PCLK = HSI = 8~MHz \cite[\S6.2.2, \S6.4.2]{rm0091}, and BR[2:0] divides it by
$2^{\mathrm{BR}+1}$ \cite[\S28.9.1]{rm0091}:
\begin{equation}
  f_\mathrm{SCK} = \frac{f_\mathrm{PCLK}}{2^{\mathrm{BR}+1}}
                 = \frac{8~\text{MHz}}{2^{\BRDec+1}}
                 = \fSCKPred~\text{kHz}.
  \label{eq:sck}
\end{equation}

\paragraph{Alternate functions.}
PB13--PB15 are set to alternate-function mode (MODER = \texttt{10})
\cite[\S8.4.1]{rm0091} and AF0 is selected in \reg{GPIOB\_AFRH}
(AFSEL13--15 = \texttt{0000}) \cite[\S8.4.10]{rm0091}, which connects them to
SPI2 \cite[Tab.~15]{stm32f051ds}. PB12 (CS) is a general-purpose output.
